\subsection{Regional comparisons and repertoire composition}
We compare 2017 with 2021 in the fixed 55-market panel. Regional summaries give each market equal weight and retain a fixed membership: Latin America (16), Nordic countries (5), East Asia (Japan and Taiwan only), Southeast Asia (5), six Anglophone markets and the 21 remaining panel markets. Membership lists are exported. These groups describe the available panel; they are not representative samples of entire world regions.

For each theme, a symmetric decomposition separates changes within repertoire groups from changes in the groups' shares. The within-group term weights the change in thematic prevalence by the group's average share in the two years. The composition term weights the change in group share by its average thematic prevalence. Their sum exactly reproduces the total change. We calculate this separately by broad genre and catalogue release group, within each market, then average markets. The two decompositions answer different questions and must not be added together.

Overlapping artist tags are assigned to one broad genre through a fixed priority rule published in the analysis script; untagged songs remain a group. Genre tags cover 55.8\% of the 199,496 classified endpoint recording--market--year observations. Catalogue release dates cover 75.3\%; dates later than the recording's first chart year are treated as unknown. Release groups are through 2016, 2017--2018, 2019--2020, 2021 onward and unknown. Catalogue dates can refer to editions or reissues rather than first artistic release. When a group appears in only one endpoint, its contribution is assigned to composition, with no invented within-group change. Actual lyric languages are absent, so no language decomposition is claimed.

Regional changes and decomposition terms have 95\% percentile intervals from 400 recording-level Poisson bootstrap replicates, using seed 20261012. Each recording receives one weight shared across all its markets and years. This preserves the dependence caused by international hits. Intervals condition on the observed markets and frozen AI labels; they do not include unmeasured classification error or establish population-level cultural change.

\subsection{Theme clusters and chart longevity}
We cluster the 68 markets using all 26 theme percentages from the combined historical period, 2017 through 13 March 2022. We standardize each dimension across markets before Ward hierarchical clustering, compare two through eight clusters and choose the largest mean silhouette score, a measure of separation. This is an exploratory summary without external or resampling-based cluster validation.

Longevity uses the classified MGD Top-200 panel in 51 markets. We sum observed daily chart appearances from 2017 through 2021 for each recording--market pair, counting re-entries. The dataset contains 224,598 pairs and 78,275 recording IDs. Themes can overlap. A linear model of log(1 + chart days) includes indicators for romance, explicitness and conspicuous consumption together, plus market, first observed chart year, broad genre and catalogue release-group effects. Standard errors group all appearances of a recording together across markets. Songs already charting at the start or still charting at the end have incomplete lifetimes; the outcome is observed chart presence, not a survival estimate or a causal effect of lyrical content.
